\chapter{Uvod}
\label{sec:uvod}

Za učenje iz podatkov, kjer niso pomembne zgolj lastnosti posameznih elementov, ampak tudi odnosi med njimi, lahko uporabimo grafovske nevronske mreže (GNN)~\cite{gilmer2017mpnn,kipf2017gcn}.
Primer takšnih podatkov so podatki sledilnika pogleda.
Meritvam sledilnika pogleda lahko smiselno določimo časovne in prostorske odnose in jim s tem priredimo časovno-prostorski graf. 
Na tem grafu se z grafovsko nevronsko mrežo učimo predstavitev vozlišč in grafov, ki jih uporabimo za klasifikacijo grafa.
Podatki, uporabljeni v tej nalogi, so označeni s čustvenimi oznakami valence, zato GNN učimo binarno klasifikacijo valence, rezultate pa primerjamo z rezultati uveljavljenih negrafovskih pristopov. 
Eksperimentalno primerjavo izvedemo na štirih signalnih naborih: samo koordinatah pogleda (pogled), samo velikostih zenic (zenici), združenih koordinatah pogleda in velikostih zenic (pogled+zenici) ter celotnem naboru razpoložljivih signalov (vsi signali).

\section{Raziskovalni problem}
\label{subsec:raziskovalni-problem}

Raziskovalni problem naloge je zasnova, implementacija in ovrednotenje grafovske predstavitve podatkov sledilnika pogleda za klasifikacijo valence.


\subsection{Raziskovalna vprašanja}
\label{subsec:raziskovalna-vprasanja}

Glavno raziskovalno vprašanje je:

\begin{enumerate}[label=\textbf{RV\arabic*:}]
    \item Kakšna je klasifikacijska uspešnost grafovske nevronske mreže za binarno klasifikacijo valence iz podatkov sledilnika pogleda v primerjavi s klasifikacijsko uspešnostjo uveljavljenih negrafovskih pristopov?
\end{enumerate}

Dodatna podvprašanja so:

\begin{enumerate}[label=\textbf{PV\arabic*:}]
    \item Kako podatke sledilnika pogleda smiselno predstaviti kot časovno-prostorski graf?
    \item Kako se uspešnost klasifikatorjev razlikuje med posameznimi nabori signalov sledilnika pogleda?
    \item Kako se klasifikacijska uspešnost spreminja z večjo kompleksnostjo grafovske nevronske mreže?
\end{enumerate}

\section{Doprinosi naloge}
\label{subsec:prispevki}

Glavni doprinosi diplomske naloge so:

\begin{enumerate}
    \item zasnova časovno-prostorske grafovske predstavitve časovnih oken podatkov sledilnika pogleda;
    \item ovrednotenje arhitekturne lestvice grafovskih modelov za klasifikacijo čustvenih oznak iz podatkov sledilnika pogleda;
    \item primerjava klasifikacijske uspešnosti pri različnih signalnih naborih;
    \item primerjava grafovskih modelov z negrafovskimi osnovnimi modeli in zamrznjenimi prednaučenimi predstavitvenimi modeli.
\end{enumerate} 

\END_NOTE{Ali rabimo tele doprinose ali raje odstranimo, glede na to da se samo ponovijo raziskovalna vprašanja?}

\section{Struktura naloge}
\label{subsec:struktura}

V poglavju~\ref{sec:ozadje} predstavimo teoretično ozadje grafov, grafovskih nevronskih mrež, podatkov sledilnika pogleda in afektivnega računalništva.
V poglavju~\ref{sec:sorodna-dela} predstavimo sorodna dela, vključno z uporabo podatkovne zbirke \MAHNOB{} in metodami učenja predstavitev iz podatkov pogleda.
Poglavje~\ref{sec:podatki} opisuje uporabljene podatke, ciljne oznake in predobdelavo.
V poglavju~\ref{sec:grafovska-predstavitev} opišemo konstrukcijo časovno-prostorskega grafa.
Poglavje~\ref{sec:modeli} predstavi uporabljene modele.
V poglavju~\ref{sec:eksperimenti} opišemo eksperimentalni protokol, v poglavju~\ref{sec:rezultati} pa rezultate glavne primerjave modelov in signalnih naborov.
Poglavje~\ref{sec:diskusija} vsebuje interpretacijo rezultatov, omejitve in možnosti nadaljnjega dela.
Nalogo sklenemo v poglavju~\ref{sec:zakljucek}.
